% TOP_PROBLEM: {"id": 3100082, "title": "Consider p(v, t), the probability that a walk beginning from the origin ends at the point v on the d-dimensional integer", "category_id": 2, "category": {"id": 2, "name": "combinatorics", "display_name": "Combinatorics", "description": "Counting problems, graph theory, discrete structures.", "slug": "combinatorics", "order_index": 2, "created_at": "2026-07-31T15:26:25.671Z"}, "status": "open", "problem_number": "AMR-030-0082", "set_id": 13}
% FIRST_POSTED: 2026-10-02
% ENTRY_KIND: statement_only
% UnsolvedMath ID 3100082; source catalogue code AMR-030-0082
% !TeX program = lualatex
% Definition and sources only; this file is not an LLM solution attempt.
\documentclass[11pt,a4paper]{article}
\usepackage[margin=25mm]{geometry}
\usepackage{fontspec}
\setmainfont{lmroman10-regular.otf}[BoldFont=lmroman10-bold.otf,
 ItalicFont=lmroman10-italic.otf,BoldItalicFont=lmroman10-bolditalic.otf]
\usepackage{amsmath,amssymb,mathtools}
\usepackage{unicode-math}
\setmathfont{latinmodern-math.otf}
\AtBeginDocument{\let\setminus\smallsetminus}
\usepackage{xurl,hyperref}
\hypersetup{colorlinks=true,urlcolor=blue}
\setcounter{secnumdepth}{0}
\setlength{\parindent}{0pt}
\setlength{\parskip}{5pt}
\setlength{\emergencystretch}{3em}
\begin{document}
\section*{Consider p(v, t), the probability that a walk beginning from the origin ends at the point v on the d-dimensional integer}\label{3100082}
{\small UnsolvedMath ID 3100082 · AMR-030-0082 · Combinatorics · AMR Open Problem Lists}
\subsection{Definitions and mathematical statement}
Let $(S_t)_{t\ge0}$ be simple symmetric nearest-neighbor random
walk on $\mathbb Z^d$, with $S_0=0$ and independent increments
uniformly distributed over $\{\pm e_1,\ldots,\pm e_d\}$.
For a fixed site $v\in\mathbb Z^d$, write
$p(v,t)=\Pr(S_t=v)$ and $m=\|v\|_1=\sum_i|v_i|$.
The probability is zero when $t<m$ or $t\not\equiv m\pmod2$.
Remove these forced parity zeros by putting
\[
                       q_j=p(v,m+2j)\qquad(j\ge0).
\]
Is this sequence unimodal for every $d\ge1$ and every site $v$?
Precisely, must there be an integer $J\ge0$ with
\[
 q_0\le q_1\le\cdots\le q_J
                  \quad\text{and}\quad
 q_J\ge q_{J+1}\ge q_{J+2}\ge\cdots\,?
\]
Equal adjacent values are permitted, so a plateau at the maximum
does not violate unimodality.

Dimension and target site remain fixed as time varies. The
question concerns the probability of being at the site at a
given time, including trajectories that visited it earlier;
it is not a first-hitting-time distribution. Deleting the
parity-forced zeros is essential, because the full sequence
usually alternates between positive values and zero. No
monotonicity in the location $v$, or in the dimension, is asserted.
\subsection{Short English statement}
After removing the forced parity zeros, is the probability that simple lattice random walk is at a fixed site unimodal in time?
\subsection{Sources}
\begin{itemize}
\item[S1] ULAM.AI and the UnsolvedMath Contributors, supplied problems.json, record 3100082 (AMR-030-0082), AMR Open Problem Lists. Catalogue identity and source transcription; CC BY 4.0.
\url{https://huggingface.co/datasets/ulamai/UnsolvedMath}.
\item[S2] Cooper - Combinatorial Problems I Like (2020 snapshot), Probability, source-order bullet 82. Original source indexed by AMR-030-0082.
\url{https://people.math.sc.edu/cooper/combprob.html}.
\end{itemize}
\end{document}
